% =========================================================
% Branded Whitepaper Template (LaTeX)
% Target: clean, professional, LinkedIn-ready PDF
% Engine: pdflatex (works) | lualatex/xelatex (optional)
% =========================================================
\documentclass[11pt]{article}

% -------------------------
% Page + typography
% -------------------------
\usepackage[letterpaper,margin=1in]{geometry}
\usepackage{microtype}
\usepackage{setspace}
\setstretch{1.12}
\usepackage[T1]{fontenc}
\usepackage[utf8]{inputenc} % omit if using lualatex/xelatex
\usepackage{lmodern}

% -------------------------
% Branding + visuals
% -------------------------
\usepackage{xcolor}
\definecolor{BrandAccent}{HTML}{204060} % deep blue-gray (edit to taste)
\definecolor{BrandGray}{HTML}{6E6E6E}
\definecolor{RuleGray}{HTML}{D9D9D9}
\usepackage{graphicx}

\newcommand{\brandlogo}{%
  % Optional: replace with \includegraphics[height=0.35in]{logo.png}
  % Keep subtle for LinkedIn readability.
  \textcolor{BrandAccent}{\textbf{WHITEPAPER SERIES}}%
}

% -------------------------
% Headers/footers
% -------------------------
\usepackage{fancyhdr}
\pagestyle{fancy}
\fancyhf{}
\renewcommand{\headrulewidth}{0pt}
\fancyhead[R]{\small\textcolor{BrandGray}{\brandlogo}}
\fancyfoot[C]{\small\textcolor{BrandGray}{\thepage}}
\fancyfoot[R]{\small\textcolor{BrandGray}{\copyright\ \AuthorName\ |\ LinkedIn Publication}}

% -------------------------
% Section styling
% -------------------------
\usepackage{titlesec}
\titleformat{\section}
  {\Large\bfseries\color{BrandAccent}}
  {}{0pt}{}
\titlespacing*{\section}{0pt}{16pt}{8pt}

\titleformat{\subsection}
  {\large\bfseries\color{BrandAccent}}
  {}{0pt}{}
\titlespacing*{\subsection}{0pt}{12pt}{6pt}

% -------------------------
% Links (optional but nice)
% -------------------------
\usepackage[hidelinks]{hyperref}
\hypersetup{
  colorlinks=true,
  linkcolor=BrandAccent,
  urlcolor=BrandAccent,
  citecolor=BrandAccent
}

% -------------------------
% Callout box
% -------------------------
\usepackage[most]{tcolorbox}
\tcbset{
  colback=white,
  colframe=BrandAccent,
  boxrule=0.8pt,
  arc=2pt,
  left=10pt,right=10pt,top=8pt,bottom=8pt
}
\newtcolorbox{keytakeaway}{}

% -------------------------
% Simple executive summary environment
% -------------------------
\newenvironment{execsummary}{
  \section*{Executive Summary}
  \addcontentsline{toc}{section}{Executive Summary}
  \vspace{-2pt}
  \color{black}
}{\vspace{6pt}}

% -------------------------
% Metadata (edit these)
% -------------------------
\newcommand{\PaperTitle}{Strategic Mastery: Writing That Moves Decisions}
\newcommand{\PaperSubtitle}{A practical playbook for producing high-impact, auditable strategy artifacts in GitHub}
\newcommand{\AuthorName}{Adam}
\newcommand{\AuthorRole}{Operator / Strategist}
\newcommand{\PaperDate}{January 2026}
\newcommand{\SeriesTag}{Whitepaper Series}
% -------------------------
% Title page
% -------------------------
\newcommand{\maketitlepage}{
  \thispagestyle{empty}
  \vspace*{0.6in}

  % Optional logo area (left) + series tag (right)
  \noindent
  \begin{minipage}[t]{0.65\textwidth}
    \textcolor{BrandGray}{\small \SeriesTag}
  \end{minipage}
  \begin{minipage}[t]{0.35\textwidth}
    \raggedleft
    \textcolor{BrandGray}{\small \brandlogo}
  \end{minipage}

  \vspace{0.35in}

  {\Huge\bfseries\color{BrandAccent}\PaperTitle\par}
  \vspace{0.12in}
  {\large\color{BrandGray}\PaperSubtitle\par}
  \vspace{0.25in}

  \noindent\textcolor{RuleGray}{\rule{\textwidth}{0.8pt}}
  \vspace{0.2in}

  {\normalsize
    \textbf{\AuthorName}\par
    \textcolor{BrandGray}{\AuthorRole}\par
    \vspace{6pt}
    \textcolor{BrandGray}{\PaperDate}\par
  }

  \vfill

  \noindent\textcolor{BrandGray}{\small
    \textbf{Distribution:} Public (LinkedIn).\quad
    \textbf{Disclaimer:} Views are the author's and do not represent official policy.
  }
  \newpage
}

% -------------------------
% Document
% -------------------------
\begin{document}

\maketitlepage

% Optional table of contents (recommended for 5+ pages)
% \tableofcontents
% \newpage

\begin{execsummary}
  A whitepaper should not merely explain. It should produce a decision.

  \begin{itemize}
    \item \textbf{Problem:} Many “strategy documents” are not falsifiable, not owned, and not operational.
    \item \textbf{Insight:} Strategy becomes real only when it is translated into constraints, measurements, and commitments.
    \item \textbf{Why it matters:} Unowned narratives create drift: slow decisions, duplicated work, and security/operational debt.
    \item \textbf{What to do:} Standardize an auditable format (GitHub + LaTeX), require testable recommendations, and instrument outcomes.
  \end{itemize}
\end{execsummary}

\section{Problem Statement}
Most “whitepapers” fail because they optimize for sounding smart instead of changing behavior. They provide context without commitments, analysis without a decision model, and recommendations without owners.

This repository is designed to produce strategy artifacts that are: (1) readable by senior leaders, (2) actionable by operators, and (3) auditable via git history.

\section{Context and Constraints}
Strategy writing that affects change must operate within real constraints:
\begin{itemize}
  \item \textbf{Time:} Leaders skim. You have minutes, not hours.
  \item \textbf{Attention:} The document competes against meetings, crises, and incentives.
  \item \textbf{Governance:} Enterprise systems move at the speed of risk management and procurement.
  \item \textbf{Accountability:} Without owners and timelines, recommendations decay into vibes.
\end{itemize}

\section{Analysis and Key Insights}
\subsection{Insight 1: Make claims falsifiable}
A strong paper makes explicit claims that could be proven wrong. This forces clarity about assumptions and makes disagreement productive.

\subsection{Insight 2: Convert narrative into a decision structure}
Every section should reduce uncertainty for a specific decision: \emph{what to do next, by whom, and how success will be measured}.

\begin{keytakeaway}
  \textbf{Key Takeaway:} If the reader cannot name the owner, timeline, and metric after reading, the paper didn't do its job.
\end{keytakeaway}

\section{Implications}
When artifacts become auditable and operational:
\begin{itemize}
  \item The organization builds a compounding library of decisions (not just opinions).
  \item Disagreements become about assumptions and data, not status and rhetoric.
  \item Execution accelerates because work packages fall directly out of recommendations.
\end{itemize}

\section{Recommendations}
Make recommendations specific, testable, and operational:
\begin{enumerate}
  \item \textbf{Adopt a standard format} for strategy briefs (this template) and require a one-page executive summary.
  \item \textbf{Require ownership:} every recommendation must have an owner, a deadline, and a success metric.
  \item \textbf{Instrument outcomes:} define what will be measured in 30/60/90 days and where it will be reported.
\end{enumerate}

\section{Conclusion}
High-impact writing is a lever: it compresses ambiguity into decisions. This repository is a starting point for producing a series of strategy documents that are designed to be read, debated, and executed.

\end{document}
